\documentclass[12pt]{article}
\usepackage[utf8]{inputenc}
\usepackage[T1]{fontenc}
\usepackage{hyperref}
\title{Kurnai Language Resources for K-12 Curriculum Development}
\author{Compiled by Claude Sonnet 5 (claude-sonnet-5) for Andr\'as Kornai}
\date{}
\begin{document}
\maketitle
\section{Introduction}
Kurnai (Gunaikurnai, Gunnai; ISO 639-3 unn) is a sleeping language of Gippsland, Victoria, in a revival phase. Wikipedia lists 4 speakers in the 2016 census. Only a few resources appeared in search, and no dictionary or wordlist was found.
\section{Primary Educational Resources}
Bairnsdale West Primary School has an overview of its Gunaikurnai language and cultural perspectives work. \url{https://bwestps.vic.edu.au/learning-and-wellbeing/gunaikurnai-language-and-culture}
\section{Linguistic and Reference Materials}
General description: \url{https://en.wikipedia.org/wiki/Gunaikurnai_language}. The Endangered Languages Project has records and a bibliography. \url{https://endangeredlanguages.com/lang/6694/bibliography?hl=pt-br}
\section{Religious and Community Resources}
The school's overview points to the GLaWAC and Bataluk Trail websites. I did not see their URLs in results, so none is given here. The Kurnai Language Program (Yirruk-Tinnor) has run for over 20 years, with a page hosted by the Victorian Aboriginal Corporation for Languages. I did not see that page's URL either.
\section{Digital and Audio Resources}
The Endangered Languages Project lists a video report on the revival program and the story \emph{Borun \& Tuk: A Gunnai Dreaming Story}. \url{https://endangeredlanguages.com/node/40129}
\section{Curriculum Development Context}
Materials are scarce. Any curriculum should be built with the Gunaikurnai community, through GLaWAC and the Kurnai Language Program.
\end{document}